\section{Loops on a sphere and on a torus}\label{sec:loop}
A torus carries a loop that goes around its first circle once. On a two-dimensional sphere, every uniformly presented loop can be contracted. Both assertions have finite arithmetic proofs. Together they distinguish the two surfaces even when the hypothetical maps are given by evaluators and need not be simplicial.

\subsection{Counting turns by rational differences}
Write $C=\mathbb R/\mathbb Z$ using the dense rational presentation $\mathbb Q/\mathbb Z$. At rational inputs set
$d_C(x,y)=\min(|x-y|,1-|x-y|)$ for representatives in $[0,1)$. Completion uses the metric conventions of Section~\ref{sec:metric-names}. A loop is a uniformly presented map on $[0,1]$ with equal endpoint names. A free homotopy of loops is a uniformly presented map on the parameter square with equal values on the two vertical sides at each height. Boundary equalities are equalities to every accuracy. The unit loop is $t\mapsto t\bmod 1$.

A perimeter parameter identifies $C$ with the boundary of $[0,1]^2$: traverse its four unit edges at constant speed four, starting at the lower left corner. This map is $4$-Lipschitz into the sup metric. For the inverse, two points on the same edge are joined there. On adjacent edges the route through their common corner has length at most twice their sup distance. On opposite edges that distance is one and a shorter boundary route has length at most two. Dividing the length by four bounds their circular distance. The rational maps are inverses, and these bounds extend them to names, including corners. The radial homeomorphism in the metric section identifies this square with the round circle after centering and rescaling. Thus $C^2$ (here the Cartesian square of the circle) gives the usual product presentation of the torus.

\begin{lemma}[A winding obstruction]\label{lem:winding}\label{idx-surfaces10}
The unit loop in $C$ admits no uniformly presented free contraction.
\end{lemma}
\begin{proof}
Suppose a contraction is given. From its square-uniform modulus choose a square mesh $1/N$ such that adjacent true images differ by at most $1/16$, with $N\ge 16$. At each mesh point use one rational circle reading accurate to $1/32$. On the right repeat the left reading. On the bottom use $i/N\bmod 1$ exactly. On the top repeat one reading of the constant endpoint. All replacements are still accurate to $1/32$, by the side and endpoint contracts. Along every oriented mesh edge $pq$ there is a unique rational difference $\delta(p,q)$ congruent modulo one to $q-p$ and of absolute value $<1/4$: the circular distance of its readings is at most $1/8$. Compute it by testing three lifts, with at most one admissible. Reversing an edge negates its difference. The sum around a cell is an integer because the differences telescope modulo one. Its absolute value is $<1$, so it vanishes. Sum over all cells. Interior edges cancel and the left and right sums cancel because the readings coincide. The bottom sum in the positive horizontal direction is $N(1/N)=1$, including the wrap across zero. The top sum is zero. This contradicts the equality of the two sums.
\end{proof}

\paragraph{Other degrees.}
For the loop $t\mapsto kt\bmod 1$ with $k\in\mathbb Z$, take $N>8(1+|k|)$ in this argument. Each bottom difference is $k/N$, even at a wrap. Any two such loops joined by a uniformly presented free homotopy consequently have equal degrees. Conversely equal integers give the same loop. More generally every uniformly presented loop is freely homotopic to one of these, with a computable degree. Sample rational circle readings $r_i$ at mesh vertices with repeat at the seam, choosing mesh oscillation and reading error at most $1/32$ each. Lift $r_0$ to $a_0\in[0,1)$ and cumulatively add the short edge differences to get $a_i$. Then $k=a_N-a_0\in\mathbb Z$. Interpolate the $a_i$ linearly to $A(t)$. Throughout each mesh interval the circular distance from $A(t)\bmod1$ to the original loop is at most $3/32+2/32<1/4$. Two such close points have a unique signed short displacement. It is evaluated on names by approximating to any demanded accuracy and reducing their difference to the nearest lift; the strict $1/4$ margin avoids ambiguity at a half-integer. Multiplying that displacement by the homotopy parameter joins the two loops. Next interpolate linearly in the covering line from $A(t)$ to $kt$, then reduce modulo one. At $t=1$ both lifts differ from their values at $t=0$ by $k$, so this is a free homotopy throughout. For loops based at $0$, use the exact first and last readings $0$, with $a_0=0$. The same construction fixes the basepoint. For concatenated loops concatenate the polygonal lifts, translating the second to start at the endpoint of the first. Endpoint differences add. Reversal negates the difference. This proves $\pi_1(C,0)\cong\mathbb Z$ in this presentation. Coordinate projection and pairing of the homotopies give $\pi_1(C^r,0)\cong\mathbb Z^r$ for each finite $r$, including $r=0$.

\subsection{Choosing a cone point by a bounded scan}
For $d\ge3$ let $S_\infty^{d-1}$ be the boundary of $[-1,1]^d$. Write $\rho(u)=u/\|u\|_\infty$ when a positive lower bound on the denominator is available.
\begin{theorem}[Contraction on the cube boundary]\label{thm:loopSphere}\label{idx-surfaces22}
Every uniformly presented loop on $S_\infty^{d-1}$, $d\ge3$, has a uniformly presented free contraction, primitive recursive from its fixed evaluation and modulus routines. For based loops a contraction fixing the basepoint can also be produced.
\end{theorem}
\begin{proof}
Choose $m$ so the oscillation of $\gamma$ on a mesh interval is at most $1/16$. Read rational vertices $p_i$ within $1/16$ of $\gamma(i/m)$, repeating $p_0$ at $i=m$. They may be chosen on the cube boundary; the argument also works for readings in the ambient cube with the indicated error. For the polygonal interpolant $P$, each endpoint of its edge is within $1/8$ of the true value at any point of that interval. The convexity of the sup ball gives $\|P(t)-\gamma(t)\|_\infty\le1/8$. Thus $\rho((1-s)\gamma(t)+sP(t))$ joins $\gamma$ to $\rho(P)$ with denominator at least $7/8$.

For each of the $m$ edges find a nonzero rational covector $n_i$ perpendicular to both endpoints, supported on the first three coordinates. Take the cross product of the first three coordinates; if it vanishes, choose a coordinate vector perpendicularized by crossing it with a nonzero one of the two triples. If both triples vanish, use $(1,0,0)$. Exact rational equality decides all branches. For an integer $j$ put $u(j)=(1,j,j^2,0,\dots,0)$. The nonzero polynomial $n_i\cdot u(j)$ has at most two distinct integer roots: subtract two root equations, factor the nonzero differences, and a third root would force all three coefficients to vanish. Hence some $0\le j\le 2m$ avoids the roots for every edge. Set $w=u(j)/\|u(j)\|_\infty$ for the first such $j$.

The second homotopy is $\rho((1-s)P(t)+sw)$. Here is an explicit denominator bound on edge $i$. For $s\le1/8$, $\|P(t)\|_\infty\ge7/8$ and
\[
 \|(1-s)P(t)+sw\|_\infty\ \ge (7/8)(7/8)-1/8>1/2.
\]
For $s\ge1/8$, pair with $n_i$. The pairing with $P$ vanishes, giving the lower bound $|n_i\cdot w|/(8\|n_i\|_1)$. The minimum of these finitely many positive rationals and $1/2$ is an explicit positive $b$. The bound extends to all real parameters by rational approximation, with no decision about a real comparison required. On vectors of norm at least $b$, radial normalization is $(2/b)$-Lipschitz: subtract $u/\|u\|-v/\|u\|$, then bound the remaining term by $|\|v\|-\|u\||/\|u\|$. The polygon has a computable Lipschitz bound $m\max_i\|p_{i+1}-p_i\|_\infty$; ordinary product and sum estimates now give moduli for both homotopies. They agree at the join, and all stages are periodic in $t$.

For clarity, an explicit conversion to a based homotopy uses the track $h(s)=H(0,s)$ of the free contraction. At stage $s$ traverse $h|_{[0,s]}$, the loop $H(-,s)$, and the reverse of that track, on three equal parameter intervals. At $s=0$ shorten the two constant waiting intervals by a linear change of the time parameter to recover the original loop. At $s=1$ the middle loop is constant. Retract the outgoing and returning parts simultaneously by replacing their final track parameter $1$ by $1-s$, with a constant middle third at $h(1-s)$. This fixes both endpoints and terminates at the constant base loop. Every reparameterization has a rational piecewise formula and a modulus from the input.
\end{proof}

\begin{corollary}[Sphere and torus]\label{cor:sphereTorus}\label{idx-surfaces39}
The presented sphere $S^n$, $n\ge2$, has trivial based fundamental group. It is not uniformly presented homotopy equivalent to the torus $C^2$. In particular there is no homeomorphism with equipped uniform evaluators in both directions. The same distinction applies to $C^r$ for $r\ge1$.
\end{corollary}
\begin{proof}
Use radial normalization between round and cubical boundaries. If $f:C^r\to S^n$ and $g:S^n\to C^r$ had $gf$ uniformly homotopic to the identity, apply the sphere contraction to $f$ of the loop $(t,0,\dots,0)$. Compose it with $g$, attach the given homotopy at the first end, then project to the first circle. The result would contract the unit loop, contradicting Lemma~\ref{lem:winding}. The proof needs only this one composite homotopy.
\end{proof}

\paragraph{First-order interpretation.}
A construction yielding a homotopy is a recipe for PR evaluation and modulus \emph{terms} given fixed terms for the input. In the winding argument take parameters coding just the polygon, the endpoint witnesses and the rational grid for the allegedly contracting map. The test on every face is a bounded scan. Its conclusion is an equality of integer numerators after a denominator product. For the sphere contraction the exceptional value search has exactly $2m+1$ candidates; the count follows by summing the bit for each pair $(i,j)$ in either order. Bounds on $1/b$ are computable from the obtained positive rational. This explains why the proof can choose a pole depending on the loop without any principle selecting a point outside an arbitrary countable set.
